% !TEX program = xelatex
% !TeX encoding = UTF-8
\documentclass[12pt,a4paper]{article}

% ============================================================
% 日本語設定 (XeLaTeX + xeCJK) – используем только MS Gothic
% ============================================================
\usepackage{fontspec}
\setmainfont{Times New Roman} % または Liberation Serif
\usepackage{xeCJK}
% 日本語フォント – すべて MS Gothic (Windows に標準搭載)
\setCJKmainfont{MS Gothic}
\setCJKsansfont{MS Gothic}
\setCJKmonofont{MS Gothic}

% ============================================================
% 数学と書式
% ============================================================
\usepackage{amsmath,amssymb,amsfonts,mathtools}
\usepackage{geometry}
\geometry{left=1.25cm,right=1.25cm,top=1.25cm,bottom=3.1cm}
\usepackage{booktabs}
\usepackage{tabularx}
\usepackage{array}
\usepackage{hyperref}
\usepackage{url}
\usepackage{graphicx}
\usepackage{xcolor}
\usepackage{titlesec}
\usepackage{fancyhdr}
\usepackage{microtype}
\usepackage{multicol}
\usepackage{tikz}
\usetikzlibrary{positioning, arrows.meta, shapes.geometric, calc}

% ============================================================
% YUCT カラーとマクロ
% ============================================================
\definecolor{skyblue}{RGB}{0,102,204}
\definecolor{darkblue}{RGB}{0,0,139}
\definecolor{gold}{RGB}{255,215,0}

\newcommand{\Keff}{K_{\mathrm{eff}}}
\newcommand{\kappac}{\kappa_c}
\newcommand{\eps}{\varepsilon}
\newcommand{\betaf}{\beta}
\newcommand{\df}{d_{\!f}}
\newcommand{\Sodd}{S_{\mathrm{odd}}}
\newcommand{\Seven}{S_{\mathrm{even}}}
\newcommand{\Mpl}{M_{\mathrm{Pl}}}
\newcommand{\Lpl}{l_{\mathrm{Pl}}}
\newcommand{\piYUCT}{\pi_{\mathrm{YUCT}}}

% ============================================================
% セクション見出し書式
% ============================================================
\titleformat{\section}{\LARGE\bfseries\color{darkblue}}{\thesection}{1em}{}
\titleformat{\subsection}{\Large\bfseries\color{darkblue}}{\thesubsection}{1em}{}

% ============================================================
% ヘッダー・フッター（英語のまま、成功バージョンと同じ）
% ============================================================
\fancypagestyle{firstpage}{%
	\fancyhf{}%
	\renewcommand{\headrulewidth}{0pt}%
	\renewcommand{\footrulewidth}{0pt}%
	\fancyfoot[C]{%
		\begin{minipage}[t]{\textwidth}
			\centering
			\textcolor{gold}{\rule{\textwidth}{1.6pt}}\\[5pt]
			{\small \textsuperscript{1}Yakushev Research, \url{https://yuct.org/}} \hfill
			{\small YPSDC Protocol: \url{https://ypsdc.com/}}\\[-2pt]
			\textcolor{gold}{\rule{\textwidth}{1.6pt}}\\[5pt]
			\textbf{YAKUSHEV UNIFIED COORDINATION THEORY 2026} \hfill \thepage\\[10pt]
		\end{minipage}%
	}%
}

\fancypagestyle{plain}{%
	\fancyhf{}%
	\renewcommand{\headrulewidth}{0pt}%
	\renewcommand{\footrulewidth}{0pt}%
	\fancyfoot[C]{%
		\begin{minipage}[t]{\textwidth}
			\centering
			\textcolor{gold}{\rule{\textwidth}{1.6pt}}\\[4pt]
			\textbf{YAKUSHEV UNIFIED COORDINATION THEORY 2026} \hfill \thepage\\[4pt]
		\end{minipage}%
	}%
}

\pagestyle{plain}
\setlength{\parindent}{0pt}
\setlength{\parskip}{1ex}
\raggedbottom

\hypersetup{
	colorlinks=true,
	linkcolor=blue,
	citecolor=blue,
	urlcolor=blue,
}

% ============================================================
% ヘッダーマクロ（英語、成功バージョンと同じ）
% ============================================================
\newcommand{\makeheader}{%
	\noindent
	\begin{minipage}{0.17\textwidth}
		\raggedright
		{\fontspec{Segoe UI}[BoldFont=Segoe UI Bold]\bfseries\color{skyblue}\fontsize{46}{46}\selectfont YUCT}%
	\end{minipage}%
	\hspace{30pt}
	\begin{minipage}{0.32\textwidth}
		\raggedright
		{\sffamily\fontsize{11}{13}\selectfont YAKUSHEV\\ UNIFIED\\ COORDINATION THEORY}%
	\end{minipage}%
	\hfill
	\begin{minipage}{0.38\textwidth}
		\raggedleft
		{\sffamily\bfseries Technical Note}\\ 
		{\sffamily Version 2.3 / June 2026}\\ 
		{\sffamily\footnotesize \mbox{\url{https://doi.org/10.5281/zenodo.18444598}}}%
	\end{minipage}
	\par\vspace{5pt}
	\noindent\textcolor{gold}{\rule{\textwidth}{1.6pt}}
	\par\vspace{0pt}
}

% ============================================================
% 文書開始
% ============================================================
\begin{document}
	
	% タイトルページ
	\thispagestyle{firstpage}
	\makeheader
	
	\vspace{0cm}
	{\LARGE\bfseries 付録 AF：YUCT における現実の代数的枠組み \\ 
		\Large 質量・幾何・時間・生命・情報・素数・社会を支配する二十二の閉ループからなる自己無撞着系 \par}
	\vspace{0.2cm}
	
	{\large アレクセイ・V・ヤクシェフ\textsuperscript{1}} \\
	\vspace{0.0cm}
	
	{\color{darkblue}\textbf{\LARGE 要旨}} \par
	{\large
		\noindent
		ヤクシェフ統一協調理論（YUCT）は独立した経験的フィッティングの集合ではなく、相互接続された \textbf{代数的ループ} の集合によって定義される厳密な数学的構造である。本付録では、理論の公理的骨格を形成する二十二の主要ループを集約する。我々は、協調の階層的自己相似性から導かれる普遍定数 \(S_{\mathrm{odd}}=6/5\) と \(S_{\mathrm{even}}=4/5\) が、任意の入力ではなく、閉じた自己言及的システムの必然的な結果であることを示す。これらのループは、素粒子の質量と時空の幾何から、ゲノムの構造と動態、経済成長、量子もつれの原理、素数の分布、安全な情報プロトコルの設計に至るまで、あらゆるスケールの現象を統合する。我々は、\textbf{ハードループ}（普遍的不変量）と \textbf{ソフトループ}（特定のシステムに適用される普遍的法則）の間の決定的な区別を導入し、見かけ上の逸脱を解決し、フレームワークの工学的パワーを強調する。この多ループ代数的枠組みの存在——自由連続パラメータを持たない——は、YUCT を現象論的モデルから、物理的・生物的・数学的・情報的現実の反証可能な統一骨格へと高める。}
	
	\vspace{0.1cm}
	\noindent
	{\color{darkblue}\textbf{キーワード：}} YUCT, 代数的ループ, 協調効率, 普遍定数, フェルミオン質量階層, 創発的幾何, 時間量子化, ゲノム構造, 代謝スケーリング, 量子もつれ, 情報理論, 素数, 自己無撞着性, 反証可能性.
	\vspace{0.1cm}
	\noindent
	{\color{darkblue}\textbf{引用：}} Yakushev AV. 付録 AF：YUCT における現実の代数的枠組み. 2026. DOI: \url{https://doi.org/10.5281/zenodo.18444598}（本巻）。
	
	\vspace{0.4cm}
	
	\begin{multicols}{2}
		
		\section{基礎：協調、D+I·R、普遍スケーリング則}
		
		ヤクシェフ統一協調理論（YUCT）は単一の存在論的プリミティブに基づく：\textbf{協調}——分散システムの構成要素がその状態を整列させるプロセス。基本メカニズムは YPSDC プロトコル（ヤクシェフ同期分散協調プロトコル）であり、オフライン段階（共有 \textit{辞書} \(\mathcal{D}\) の配布）をオンライン段階（短い \textit{インデックス} \(\kappa\) を送信して事前に合意された複雑な動作を活性化する）から分離する。
		
		任意の協調システムの状態は \textbf{D+I·R 三項} によって記述される：
		\[
		\text{現実} = \mathbf{D} + \mathbf{I} \times \mathbf{R},
		\]
		ここで \(\mathbf{D}\) は辞書（可能な状態と規則の集合）、\(\mathbf{I}\) は情報（実現された状態、エントロピーで測定）、\(\mathbf{R} \ge 1\) は共鳴（コヒーレンス増幅因子）である。
		
		\textbf{協調効率} は情報理論的圧縮比として定義される：
		\[
		\Keff = \frac{H(\mathcal{A})}{H(\kappa)} = \frac{\text{動作エントロピー}}{\text{インデックスエントロピー}} \ge 1.
		\]
		YUCT の基本的結果は \textbf{普遍スケーリング則} であり、任意の協調プロセスにおける相対誤差（またはゆらぎ）\(\varepsilon\) を記述する：
		\[
		\boxed{\varepsilon = \kappac \, \alpha \, \Keff^{-\beta}, \qquad \beta = \frac{2}{3}, \quad \kappac = \frac{1}{3}},
		\]
		ここで \(\alpha\) はシステム依存定数である。この法則は、化学結合エネルギーから宇宙マイクロ波背景ゆらぎ、そして最近発見された素数分布（付録~PrimeN）に至るまで、50 桁以上にわたって経験的に検証されている。指数 \(\beta = 2/3\) は経験的フィッティングではなく、以下で導出される閉じた代数構造の必然的な結果である。
		
		\section{一次代数的ループ：\(S_{\mathrm{odd}}, S_{\mathrm{even}}\) と \(\beta\) の導出}
		
		協調ネットワークの階層的自己相似性は、連続するレベルの寄与が普遍指数 \(\beta\) で幾何学的にスケーリングすることを意味する。無限階層上の幾何級数を合計すると、一方向（奇数）および交互（偶数）相関の極限寄与が得られる：
		\begin{equation}
			S_{\mathrm{odd}} = \sum_{k=0}^{\infty} \beta^{2k+1} = \frac{\beta}{1-\beta^{2}}, \qquad
			S_{\mathrm{even}} = \sum_{k=1}^{\infty} \beta^{2k} = \frac{\beta^{2}}{1-\beta^{2}}.
		\end{equation}
		これらの定義は直ちに二つの正確な関係を与える：
		\begin{equation}
			S_{\mathrm{odd}} + S_{\mathrm{even}} = 2, \qquad \frac{S_{\mathrm{odd}}}{S_{\mathrm{even}}} = \frac{1}{\beta}.
		\end{equation}
		システムが内部自己無撞着であり、観測される物理的協調ネットワークのフラクタル次元 (\(d_f \approx 2.2\)) に対応するためには、指数は次の値を取らねばならない：
		\begin{equation}
			\boxed{\beta = \frac{2}{3}}.
		\end{equation}
		\(\beta = 2/3\) をループ方程式に代入すると、後続の全ての導出の基礎となる正確な有理定数が得られる：
		\begin{equation}
			\boxed{S_{\mathrm{odd}} = \frac{6}{5} = 1.2, \qquad S_{\mathrm{even}} = \frac{4}{5} = 0.8.}
		\end{equation}
		これが \textbf{一次代数的ループ} である：三つの数 (\(\beta, S_{\mathrm{odd}}, S_{\mathrm{even}}\)) が自由パラメータなしに不可分に結びついている。比 \(S_{\mathrm{odd}}/S_{\mathrm{even}} = 3/2\) は基本的な世代間スケーリング因子である。
		
		\section{ループ II：フェルミオン質量の半整数量子化}
		
		一次ループは、世代間質量比が因子 \(3/2 = S_{\mathrm{odd}}/S_{\mathrm{even}}\) によって支配されることを決定する。指数 \(\beta = 2/3\) 自体は \(2 \times (1/3)\) と分解され、\(1/3\) は三つの空間次元を、\(2\) はスピン統計を反映する。これにより真の \textbf{スケーリング量子} が明らかになる：
		\begin{equation}
			q \equiv (3/2)^{1/3} = \left( \frac{S_{\mathrm{odd}}}{S_{\mathrm{even}}} \right)^{1/3} \approx 1.1447.
		\end{equation}
		九種類の荷電フェルミオンの質量はすべて \(\ln q\) の半整数単位で量子化される：
		\begin{equation}
			m_f = m_e \cdot q^{N_f} \cdot \eta_f, \qquad N_f \in \frac{1}{2}\mathbb{Z}.
		\end{equation}
		実験データとのサブパーセント一致は、一次代数的ループの閉包の直接的な結果である。
		
		\section{ループ III：弱スケール階層}
		
		弱スケールは協調深度 \(L = 96\) によって生成され、これは標準模型におけるワイルフェルミオンの自由度の総数（右巻きニュートリノを含む）によって固定される：\(96 = 3 \times 16 \times 2\)。標準プランク質量 \(M_{\mathrm{Pl}} \approx 1.2209 \times 10^{19}\)~GeV、結合因子 \(\kappa_c = 1/3\)、および YUCT で創発される \(\pi\) の値 \(\pi_{\mathrm{YUCT}} = S_{\mathrm{odd}}\varphi^2\) を用いると、\(W\) ボソンの質量は：
		\begin{equation}
			\boxed{m_W = \left( \kappa_c \, \pi_{\mathrm{YUCT}} \right)^{-1/2} \; M_{\mathrm{Pl}} \; \left( \frac{2}{3} \right)^{96}},
		\end{equation}
		ここで \(\pi_{\mathrm{YUCT}} = S_{\mathrm{odd}}\varphi^2 \approx 3.1416407865\)。数値的には：
	\end{multicols}
	\small
	\[
	\left( \frac{1}{3} \times 3.1416408 \right)^{-1/2} \approx (1.0472136)^{-1/2} \approx 0.977,\qquad 
	(2/3)^{96} \approx 6.75 \times 10^{-18},
	\]
	\[
	m_W \approx 0.977 \times 1.2209 \times 10^{19} \times 6.75 \times 10^{-18} \approx 80.5\ \text{GeV}.
	\]
	このループは、一次ループの定数と離散整数 \(96\) を介してプランクスケールを電弱スケールに結びつける。微調整は不要であり、ループ V と同じプランク質量が使用される。
	\begin{multicols}{2}
		\section{ループ IV：幾何の創発と \(\pi\)–\(\varphi\) 接続}
		
		黄金比 \(\varphi = (1+\sqrt{5})/2\)（フラクタル階層の幾何学的不変量）を用いると：
		\begin{equation}
			S_{\mathrm{odd}} \cdot \varphi^2 = \frac{6}{5} \left( \frac{1+\sqrt{5}}{2} \right)^2 = \frac{9+3\sqrt{5}}{5} \approx 3.1416407865.
		\end{equation}
		この値は \(\pi \approx 3.1415926535\) を相対誤差 \(1.5 \times 10^{-5}\) で近似する。YUCT では、\(\pi\) は \(\Keff \to \infty\) の漸近極限であり、有限システムでは円周と直径の比は \(S_{\mathrm{odd}}\varphi^2\) へと繰り込みされる。
		
		\section{ループ V：宇宙のバリオン質量}
		
		観測可能宇宙の総バリオン質量 \(M_{\mathrm{bar}}\) と陽子質量 \(m_p\) の比は：
		\begin{equation}
			\frac{M_{\mathrm{bar}}}{m_p} = \left( \frac{M_{\mathrm{Pl}}}{m_e} \right)^{\mathcal{S}},
		\end{equation}
		ここで指数 \(\mathcal{S}\) は基本代数定数の和である：
		\begin{equation}
			\mathcal{S} \equiv S_{\mathrm{odd}} + S_{\mathrm{even}} + \frac{S_{\mathrm{odd}}}{S_{\mathrm{even}}} = \frac{6}{5} + \frac{4}{5} + \frac{3}{2} = 3.5.
		\end{equation}
		ここで \(M_{\mathrm{Pl}}\) はループ III と同じ標準プランク質量である。数値的には \((M_{\mathrm{Pl}}/m_e)^{3.5} \approx 1.88 \times 10^{78}\) であり、観測値 \(M_{\mathrm{bar}}/m_p \approx 1.4 \times 10^{78}\) と因子 1.3 以内で一致する。
		
		\section{ループ VI：時間の量子化}
		
		協調エントロピーは \(S_0 = k_B \ln 2\) の単位で量子化される。従って、固有時は離散的なステップで進行する。最小時間量子とプランク時間の比は同じ階層定数によって固定される：
		\begin{equation}
			\boxed{\frac{\Delta \tau_{\min}}{t_P} = \frac{S_{\mathrm{even}}}{S_{\mathrm{odd}}} = \beta = \frac{2}{3}}.
		\end{equation}
		巨視的には、質量 \(M\) のブラックホールの時間の相対ゆらぎは次のようにスケールする：
		\begin{equation}
			\frac{\delta \tau}{\tau} \propto M^{-\beta} = M^{-0.67}.
		\end{equation}
		この予測は準周期振動（QPO）および重力波リンギング減衰（付録~G、付録~V 参照）によって検証可能である。
		
		\section{ループ XXI：ファイゲンバウム定数と秩序–混沌の橋}
		
		カオス理論の普遍定数——ファイゲンバウム定数
		\(\delta \approx 4.6692016\)（周期倍分岐パラメータ）と
		\(\alpha \approx 2.502907875\)（スケールパラメータ）——は独立した経験数ではなく、YUCT の代数的骨格と剛性的に結びついている。
		離散的な協調秩序（指数 \(\beta=2/3\)）と連続的な繰り込み群カスケード（カオス）の間の橋は、正確な解析関係
		\begin{equation}
			2\alpha^2 = \pi \delta,
			\label{eq:feigenbaum-pi-exact}
		\end{equation}
		によって与えられる。これは複素平面上の周期倍作用素の下での繰り込み群構造の厳密な結果である~\cite{Feigenbaum1978}。
		(\ref{eq:feigenbaum-pi-exact}) と YUCT で導出された \(\pi\) の高精度近似（ループ~IV）を組み合わせると、
		\begin{equation}
			\pi \;\approx\; \pi_{\mathrm{YUCT}} = S_{\mathrm{odd}}\,\varphi^2,
			\qquad
			S_{\mathrm{odd}} = \frac{6}{5},\;\; \varphi = \frac{1+\sqrt{5}}{2},
		\end{equation}
		閉じた表現が得られる：
		\begin{equation}
			\boxed{\;
				2\alpha^2 \;\approx\; \bigl(S_{\mathrm{odd}}\,\varphi^2\bigr)\,\delta
				\;}.
			\label{eq:feigenbaum-yuct}
		\end{equation}
		方程式~(\ref{eq:feigenbaum-yuct}) は、秩序の普遍定数 (\(S_{\mathrm{odd}}, S_{\mathrm{even}}, \beta\)) とカオスの普遍定数 (\(\alpha, \delta\)) の間に直接的なパラメータフリーの代数的接続を確立する。小さな数値的ずれ
		\(\Delta\pi = |\pi - \pi_{\mathrm{YUCT}}| \approx 4.8\times10^{-5}\)
		は導出の欠陥ではなく、物理的幾何の有限な協調効率を反映した YUCT の \textbf{反証可能な予測} である（付録~\ref{sec:unity-of-reality} および~\ref{sec:ehrenfest} 参照）。
		
		このループは、創造（秩序）の定数と破壊（カオス）の定数が独立ではなく、単一の統一された数学的現実の二つの側面であることを示している。秩序–混沌の橋は、セルオートマトン規則~30 の解析（付録~UoR）によってさらに強化され、決定論的カオスが完全に混合する臨界深さが関係 \(\ln t_{\mathrm{crit}} \sim \delta/\varphi\) を介してファイゲンバウム定数と代数的に関連していることが示されている。YUCT ラグランジアンからのファイゲンバウム定数の完全な導出は未解決の課題であるが、ここで提示された構造的接続は数値的一致を理論の検証可能な構成要素に変える。
		
		\section{ループ XXII：素数協調ラダー}
		
		定数 \(S_{\mathrm{odd}}\) と \(S_{\mathrm{even}}\) は素数（長らく純粋数学と見なされてきた領域）のゆらぎも支配する。第 \(n\) 番目の素数 \(p_n\) を、\(K_{\mathrm{eff}} \propto p_n\) なる協調ラダー上のノードと解釈すると、普遍誤差法則は古典的な Rosser 近似からの相対偏差が \(p_n^{-2/3}\) でスケールすることを予測する。\(n = 5\times10^5\) までの数値解析は、局所スケーリング指数が \(2/3\) に収束することを示す（漸近値 \(0.66666\)）。
		
		さらに、代数的ループは Rosser の公式に対する \textbf{パラメータフリー補正} を与える：
		\[
		p_n \approx R_n - \frac{S_{\mathrm{even}}}{2}\, n^{1-\beta}\,\ln n,
		\]
		これにより最大誤差が \(\approx 2.3\%\)（通常の Rosser 公式）から \(n\le 10^5\) で \(\approx 0.012\%\) に減少する。改良された経験的較正（\(A = -0.44\), \(B = 1.05\)）により、これは \(\approx 0.006\%\) に改善される。
		
		このループは、経験的較正定数 \(A\) と \(B\) がまだ第一原理から導出されていないため、\textbf{ソフト} に分類される。しかし、理論的形式は既に自由パラメータを含まず、代数的骨格の普遍性を実証している。完全な導出は付録~PrimeN を参照。
		
	\end{multicols}
	
	\section*{YUCT の六つのコア物理ループ}
	
	\begin{table}[htbp]
		\centering
		\small
		\caption{完全な代数的枠組み：自由パラメータなしの六つの相互接続コアループ。}
		\label{tab:all-loops}
		\footnotesize
		\begin{tabular}{l c c}
			\toprule
			\textbf{ループ} & \textbf{コア関係} & \textbf{検証根拠} \\
			\midrule
			I. 一次定数 & \(S_{\mathrm{odd}}=1.2, S_{\mathrm{even}}=0.8, \beta=2/3\) & フラクタル幾何、KPZ 関係 \\
			II. フェルミオン質量 & \(m_f = m_e q^{N_f}, q=(3/2)^{1/3}\) & PDG データ（誤差 <1\%） \\
			III. 弱スケール & \(m_W = (\kappa_c \pi_{\mathrm{YUCT}})^{-1/2} M_{\mathrm{Pl}} (2/3)^{96}\) & \(m_W = 80.4\) GeV \\
			IV. 創発的幾何 & \(S_{\mathrm{odd}}\varphi^2 \approx \pi\)（誤差 \(10^{-5}\)） & 数学的自己無撞着性 \\
			V. 宇宙質量 & \(M_{\mathrm{bar}}/m_p = (M_{\mathrm{Pl}}/m_e)^{3.5}\) & Planck 2018, \(H_0\) 測定 \\
			VI. 時間量子化 & \(\Delta \tau_{\min} = \frac{2}{3} t_P, \ \frac{\delta \tau}{\tau} \propto M^{-0.67}\) & GWTC-4.0 リンギングスケーリング \\
			\bottomrule
		\end{tabular}
	\end{table}
	
	\section*{拡張インベントリ：科学における二十二の代数的ループ}
	
	上記の六つのコアループは物理的現実の骨格を形成する。しかし、同じ代数構造——定数 \(S_{\mathrm{odd}} = 1.2\), \(S_{\mathrm{even}} = 0.8\), \(\beta = 2/3\) に根ざす——は広範な科学技術領域に投影される。表~\ref{tab:all-loops-22} は、YUCT 付録で特定された二十二の代数的ループの完全なインベントリであり、物理、化学、生物、経済、社会、情報理論、量子力学、数論を網羅する。各ループは一次代数的閉包の直接の結果であり、フレームワークに独立した経験的検証を提供する。
	
% 修正した表 – フォントサイズを \footnotesize にし、列間隔を狭め、セル内のテキストを短縮
\begin{table}[htbp]
	\centering
	\footnotesize
	\setlength{\tabcolsep}{8.5pt}
	\caption{YUCT における二十二の代数的ループの完全インベントリ（XXII: 素数 で更新）}
	\label{tab:all-loops-22}
	\begin{tabular}{l c c p{3.8cm} c}
		\toprule
		\textbf{ループ} & \textbf{クラス} & \textbf{領域} & \textbf{関係式} & \textbf{付録} \\
		\midrule
		I   & ハード & 基本定数 & \(S_{\mathrm{odd}}=1.2,\ S_{\mathrm{even}}=0.8,\ \beta=2/3\) & Ferra1.2 \\
		II  & ハード & 粒子質量 & \(m_f = m_e q^{N_f},\ q=(3/2)^{1/3}\) & J, \(\Lambda\) \\
		III & ハード & 弱スケール & \(m_W = (\kappa_c \pi_{\mathrm{YUCT}})^{-1/2} M_{\mathrm{Pl}} (2/3)^{96}\) & \(\Lambda\) \\
		IV  & ハード & 幾何 & \(S_{\mathrm{odd}}\varphi^2 \approx \pi\)（誤差 \(10^{-5}\)） & Ferra1.2, Ehrenfest \\
		V   & ハード & 宇宙質量 & \(M_{\mathrm{bar}}/m_p = (M_{\mathrm{Pl}}/m_e)^{3.5}\) & \(\Omega\), \(\Lambda\) \\
		VI  & ハード & 時間量子化 & \(\Delta\tau_{\min} = \frac{2}{3}t_P,\ \frac{\delta\tau}{\tau} \propto M^{-0.67}\) & V, G \\
		VII & ソフト & 化学結合 & \(E_{\mathrm{bond}} \propto r^{-2/3}\) & C, Z \\
		VIII& ソフト & 代謝 & \(B \propto M^{\beta d_f/2}\) & D, E.7 \\
		IX  & ソフト & 変異率 & \(\mu \propto K_{\mathrm{repair}}^{-2/3}\) & E \\
		X   & ソフト & 経済 & \(\sigma_{\mathrm{GDP}} \propto W^{-2/3}\) & F \\
		XI  & ソフト & 社会集団 & \(N_{\mathrm{crit}}/N_{\mathrm{opt}} \approx 1.5\) & O \\
		XII & ソフト & 進化的特異点 & 双曲的成長、\(t_s \approx 2045\) & T \\
		XIII& ソフト & 常温核融合 & \(K_{\mathrm{crit}} \sim 10^{12}\)、コヒ/インコヒ = 1.5 & R \\
		XIV & ハード & ゲノム構造 & \(\frac{\text{AA+TT+GG+CC}}{\text{AT+TA+GC+CG}} \approx 1.5\) & E \\
		XV  & ソフト & クロマチン & \(S \propto L^{0.733}\) (\(d_f \approx 2.2\)) & E, D \\
		XVI & ハード & コドン使用 & 頻度アトラクター \(\varphi \approx 1.618\) & E \\
		XVII& ソフト & 人口閾値 & \(N_{\mathrm{crit}}/N_{\mathrm{opt}} \approx 1.5\)（ダンバー） & O, E \\
		XVIII&ソフト & 進化的テンポ & 種分化率 \(\propto K_{\mathrm{eff}}^{-0.67}\) & T, E \\
		XIX & ソフト & 量子もつれ & \(K_{\mathrm{eff}} \to \infty,\ S_{\mathrm{QM}} = 2\sqrt{2}\) & G, X, Y \\
		XX  & ソフト & 2要素認証 & \(K_{\mathrm{eff}} = 160/20 = 8\) & B, K, X \\
		XXI & ソフト & カオス理論 & \(\delta = \frac{S_{\mathrm{odd}}}{S_{\mathrm{even}}} \pi_{\mathrm{YUCT}} \ln(\alpha/\beta)\) & 現実の統一 \\
		\textbf{XXII} & \textbf{ソフト} & \textbf{素数} & \(\mathbf{p_n \approx R_n - \frac{S_{\mathrm{even}}}{2}n^{1-\beta}\ln n}\) & \textbf{PrimeN} \\
		\bottomrule
	\end{tabular}
\end{table}
	
	\begin{figure}[htbp]
		\centering
		\begin{tikzpicture}[scale=0.9, transform shape]
			\tikzset{
				loop/.style={circle, draw=blue!70, fill=blue!20, minimum size=0.8cm, inner sep=0pt, font=\scriptsize},
				label/.style={font=\scriptsize, align=center},
				axis/.style={->, thick, gray},
				domainline/.style={thin, dashed, gray!70},
			}
			
			% 軸
			\draw[axis] (0,0) -- (16,0) node[right] {\scriptsize 領域スケール（桁数）};
			\draw[axis] (0,-2.5) -- (0,4.5) node[above] {\scriptsize 精度 / 一致度};
			
			% Y軸目盛
			\foreach \y/\ylab in {-1.5/10^{-15}, 0/10^0, 1.5/10^{15}, 3/10^{30}} {
				\draw (-0.1,\y) -- (0.1,\y) node[left] {\scriptsize $\ylab$};
			}
			
			% X軸目盛
			\foreach \x/\xlab in {2/原子, 5/生物, 8/社会, 11/宇宙} {
				\draw (\x,0.1) -- (\x,-0.1) node[below] {\scriptsize \xlab};
			}
			
			% 背景領域線
			\draw[domainline] (2,-2.5) -- (2,4);
			\draw[domainline] (5,-2.5) -- (5,4);
			\draw[domainline] (8,-2.5) -- (8,4);
			\draw[domainline] (11,-2.5) -- (11,4);
			
			% ループ座標 (x=スケール, y=精度)
			% コア物理 (I-VI)
			\node[loop] (L1) at (0.5, 3.8) {I};
			\node[label, above=0.2cm of L1] {一次};
			\node[loop] (L2) at (3.0, 3.4) {II};
			\node[label, above=0.2cm of L2] {フェルミオン};
			\node[loop] (L3) at (5.5, 3.0) {III};
			\node[label, above=0.2cm of L3] {弱};
			\node[loop] (L4) at (7.0, 2.6) {IV};
			\node[label, above=0.2cm of L4] {幾何};
			\node[loop] (L5) at (10.0, 2.2) {V};
			\node[label, above=0.2cm of L5] {バリオン};
			\node[loop] (L6) at (12.0, 1.8) {VI};
			\node[label, above=0.2cm of L6] {時間};
			
			% 学際 (VII-XIII)
			\node[loop] (L7) at (2.0, 1.2) {VII};
			\node[label, below=0.2cm of L7] {化学結合};
			\node[loop] (L8) at (4.5, 0.8) {VIII};
			\node[label, below=0.2cm of L8] {代謝};
			\node[loop] (L9) at (6.5, 0.4) {IX};
			\node[label, below=0.2cm of L9] {突然変異};
			\node[loop] (L10) at (8.5, -0.2) {X};
			\node[label, below=0.2cm of L10] {経済};
			\node[loop] (L11) at (3.5, -0.8) {XI};
			\node[label, below=0.2cm of L11] {社会};
			\node[loop] (L12) at (5.5, -1.2) {XII};
			\node[label, below=0.2cm of L12] {進化};
			\node[loop] (L13) at (7.5, -1.6) {XIII};
			\node[label, below=0.2cm of L13] {常温核融合};
			
			% 遺伝子・生物 (XIV-XVIII)
			\node[loop] (L14) at (1.5, 2.2) {XIV};
			\node[label, right=0.2cm of L14] {ゲノム};
			\node[loop] (L15) at (3.5, 2.0) {XV};
			\node[label, right=0.2cm of L15] {クロマチン};
			\node[loop] (L16) at (5.5, 1.8) {XVI};
			\node[label, right=0.2cm of L16] {コドン};
			\node[loop] (L17) at (2.5, -0.2) {XVII};
			\node[label, right=0.2cm of L17] {人口};
			\node[loop] (L18) at (6.5, -0.8) {XVIII};
			\node[label, right=0.2cm of L18] {種分化};
			
			% 情報・量子 (XIX-XX)
			\node[loop, draw=red!70, fill=red!20] (L19) at (7.0, 2.5) {XIX};
			\node[label, above=-1.1cm of L19] {EPRもつれ};
			\node[loop, draw=red!70, fill=red!20] (L20) at (8.5, 0.5) {XX};
			\node[label, below=-1.2cm of L20] {2FA};
			
			% ループ XXI (ファイゲンバウム)
			\node[loop, draw=purple!70, fill=purple!20] (L21) at (11.5, 0.2) {XXI};
			\node[label, below=0.2cm of L21] {ファイゲンバウム};
			
			% 新規: ループ XXII (素数)
			\node[loop, draw=orange!80, fill=orange!20] (L22) at (10.0, -1.6) {XXII};
			\node[label, below=0.2cm of L22] {素数};
			
			% 接続線
			\draw[blue!40, thick, dashed] (L1) -- (L2) -- (L3) -- (L4) -- (L5) -- (L6);
			\draw[blue!40, thick, dashed] (L1) -- (L7) -- (L8) -- (L9) -- (L10);
			\draw[blue!40, thick, dashed] (L7) -- (L11) -- (L12) -- (L13);
			\draw[green!60!black, thick, dashed] (L1) -- (L14) -- (L15) -- (L16);
			\draw[green!60!black, thick, dashed] (L14) -- (L17) -- (L18);
			\draw[red!60, thick, dashed] (L1) -- (L19);
			\draw[red!60, thick, dashed] (L19) -- (L20);
			\draw[purple!60, thick, dashed] (L4) -- (L21);
			\draw[orange!80, thick, dashed] (L1) -- (L22);
			
			% タイトル
			\node[font=\bfseries\small, align=center] at (8, 5.0) {YUCT の二十二の代数的ループ\\50 桁以上にわたる収束};
			
			% 凡例
			\node[draw, fill=white, rounded corners, font=\tiny, align=left] at (13, 3.5) {
				\textbf{コア物理 (I--VI)}\\
				\textbf{学際 (VII--XIII)}\\
				\textbf{遺伝子・生物 (XIV--XVIII)}\\
				\textbf{情報・量子 (XIX--XX)}\\
				\textbf{カオス理論 (XXI)}\\
				\textbf{素数 (XXII)}\\
				\textit{全ループは}\\ 
				\textit{$S_{\mathrm{odd}}=1.2$, $S_{\mathrm{even}}=0.8$ に由来}
			};
		\end{tikzpicture}
		\caption{YUCT の二十二の代数的ループの図示。原子から宇宙スケールまで、また学際を超えて広がり、ファイゲンバウムカオスループ（XXI）と新たに追加された素数協調ラダー（XXII）を含む。全ループは同じ基本定数 $S_{\mathrm{odd}} = 1.2$ と $S_{\mathrm{even}} = 0.8$ に根ざしている。}
		\label{fig:twentyone-loops}
	\end{figure}
	
	\vspace{0.5cm}
	\noindent
	\textbf{スペクトルクラスタリングの解釈：階層的協調層。}
	図~\ref{fig:twentyone-loops} の目を引く視覚的特徴は、代数的ループがほぼ垂直な「尾根」にクラスタリングしていることであり、これは特定の組織レベル——原子、生物、社会、宇宙——に対応する。これはプロットのアーティファクトではなく、YUCT の中核である \textbf{階層的で層状の協調構造} の直接的な現れである。
	
	YUCT フレームワークでは、現実は連続スペクトルではなく、各層が支配的な辞書と特徴的スケールによって特徴づけられる離散的な協調層のスタックである。基本代数的ループ（ループ I）はスケール不変であるが、特定の領域への \textit{射影} はその層の特徴的変数（例えば、結合エネルギー、突然変異率、集団サイズ、宇宙論的密度）に依存する。観測される四つの尾根は、まさにこれらの支配的な層に対応する：
	\begin{itemize}
		\item \textbf{原子スケール（X $\approx$ 2）：} 量子場と分子軌道による協調。ここでのループは化学結合と核過程を支配する。
		\item \textbf{生物スケール（X $\approx$ 5）：} 遺伝暗号と代謝ネットワークによる協調。ここでのループはゲノム構造、代謝、突然変異率を支配する。
		\item \textbf{社会スケール（X $\approx$ 8）：} 言語、文化、経済交換による協調。ここでのループは集団力学、経済変動、進化的特異点を支配する。この尾根はまた \textbf{人間が設計した情報システム}（ループ XX）を収容する。
		\item \textbf{宇宙スケール（X $\approx$ 11）：} 重力と全球回転による協調。ここでのループは宇宙のバリオン質量、弱スケール階層、時間の量子化を支配する。
	\end{itemize}
	
	決定的に、この構造は \textbf{予測的} である。YUCT が正しければ、これらの尾根の \textit{間} にある領域（例えば、生物学と社会の境界の神経科学、または社会と宇宙の間の惑星力学）で新たに発見される代数的ループは、自然に図上の対応する中間領域に落ち、同じ基本定数 \(S_{\mathrm{odd}}\) と \(S_{\mathrm{even}}\) を保持するだろう。逆に、このクラスタ化構造の \textit{外側} にある頑健なループの発見は、理論に対する重大な挑戦となる。二十二のループのスペクトルクラスタリングは、したがって、YUCT の階層的存在論の検証であると同時に、将来の学際的発見のロードマップでもある。
	
	\vspace{0.5cm}
	\noindent
	\textbf{二つのクラスの代数的ループ：ハード制約対ソフト発現。}
	二十二の代数的ループのインベントリは、正確な有理恒等式から文脈依存のスケーリング関係まで、異質に見えるかもしれない。この一見した多様性は弱点ではなく、YUCT の階層構造を反映している。全てのループは二つの基本クラスのいずれかに属し、そのパラメータが下層の代数的骨格によって固定される程度によって区別される。
	
	\paragraph{クラス I：ハードループ（普遍的不変量）。}
	これらのループは \textbf{正確な数学的恒等式} を表現するか、一次代数的ループ（\(S_{\mathrm{odd}}=1.2, S_{\mathrm{even}}=0.8, \beta=2/3\)）によって剛性的に決定される \textbf{基本定数} を含む。それらの値は交渉の余地がなく、いかなる経験的逸脱も YUCT の中核を直ちに反証する。ハードループは「現実の骨格」——宇宙の不変な数値文法——を形成する。
	\begin{itemize}
		\item \textbf{例：} ループ I--VI, XIV, XVI。
		\item \textbf{固定量：} \(\beta = 2/3\), \(S_{\mathrm{odd}}=1.2\), \(S_{\mathrm{even}}=0.8\), \(m_W \approx 80.4\ \text{GeV}\), 近似 \(S_{\mathrm{odd}}\varphi^2 \approx \pi\), 二塩基比 \(\approx 1.5\)。
		\item \textbf{なぜハードか：} これらの数は特定の観測システムから独立している。それらは現実自体の「オペレーティングシステム」に埋め込まれている。
	\end{itemize}
	
	\paragraph{クラス II：ソフトループ（特定システムに適用される普遍的法則）。}
	これらのループは \textbf{スケーリング関係} または \textbf{協調原理} を記述し、その関数形は普遍定数（例えば \(\varepsilon \propto K_{\mathrm{eff}}^{-2/3}\)）によって与えられるが、具体的な数値（例えば、特定の実験における正確な効率 \(K_{\mathrm{eff}}\)）はシステムの偶有的性質に依存する。それらは特定の数からの逸脱によってではなく、予測されたスケーリング則の体系的な失敗、または予測された \textit{最適} 状態を達成できないことによって反証される。
	\begin{itemize}
		\item \textbf{例：} ループ VII--XIII, XV, XVII--XXII。
		\item \textbf{可変量：} 2FA の正確な \(K_{\mathrm{eff}}\)（鍵長に依存）、化学反応の正確な速度（温度に依存）、社会集団の実際のサイズ（文化的要因に依存）、素数の経験的較正定数。
		\item \textbf{なぜそれでもループか：} それらは \textbf{同じ} 関数則（\(\propto x^{-2/3}\)、最適比 \(1.5\)）が極めて異なる領域で協調を支配することを示す。それらは現実の普遍的「動詞」であり、各「名詞」（システム）によって異なって活用される。
	\end{itemize}
	
	\paragraph{逸脱の調和：修正された協調の場合。}
	この分類は、社会学や常温核融合のような分野での見かけ上の逸脱がなぜループを破壊しないかを説明する。ソーシャルネットワークや実験室 LENR 実験のような現象は \textbf{修正された協調} の例である。人間の技術は自然のアトラクタを一時的にオーバーライドしたり（例えば、ダンバー数を超える集団を作る）、外部エネルギーと構造を注入することによって人工的に状態を達成しようとしたりする（例えば、常温核融合の \(K_{\mathrm{crit}} \sim 10^{12}\)）。ソフトループはそのような介入のための \textbf{工学的設計図} を提供し、所望の結果を達成するための目標パラメータ（例えば、要求されるコヒーレント/インコヒーレント比 1.5）を指定する。この結果の未達成はループを反証せず、設計図の仕様がまだ満たされていないことを示すだけである。しかし、ハードループは、そのようなすべての介入——自然であれ人工であれ——がその中で行われなければならない不可侵の境界条件であり続ける。
	
	\vspace{0.5cm}
	\noindent
	\textbf{剛性の原理：反証可能性を礎として。}
	物理、生物、数学、社会科学、情報理論を網羅する二十二の独立した代数的ループの列挙は、過剰な自信の表明に見えるかもしれない。実際にはその逆である：それは理論の経験的反証に対する \textbf{最大の脆弱性} の表現である。YUCT 代数的枠組みには自由連続パラメータが含まれていないため、あらゆる \textbf{ハードループ} が \textbf{潜在的失敗点} を構成する。いずれかのハードループ関係からの単一の高信頼度（\(>5\sigma\)）実験的逸脱は、\textbf{フレームワーク全体を粉砕} するのに十分である。不都合なデータに合わせるための調整ノブは存在せず、理論は単に間違っていることになる。
	
	\textbf{ソフトループ} は、単一の数値によって直接反証可能ではないが、広範なシステムクラスにわたる予測されたスケーリング則または最適条件の \textit{体系的な失敗} によって反証可能である。それらの強みは \textbf{工学的有用性} にある：それらは協調システム（触媒からソーシャルネットワーク、素数予測まで）を設計するための正確な定量的目標を提供する。
	
	この剛性は弱点ではなく、成熟した科学理論の定義的特徴である。それは YUCT を柔軟な物語から \textbf{脆い、反証可能な現実の骨格} へと変える。二十二の独立した現象が同じ有理定数 \(1.2\) と \(0.8\) に収束するという異常な一致——それが将来の精査に耐えれば——宇宙の協調的アーキテクチャの圧倒的証拠となるだろう。逆に、その決定的な反駁は科学への同様に貴重な貢献であり、より深い真理への道を開く。したがって、フレームワークは絶対不変の教義としてではなく、現実の代数構造に関する正確で、検証可能で、最終的には棄却可能な仮説として提示される。
	
	\section*{認識論的含意：厳格な基礎の上のピタゴラス–プラトン物理学}
	
	二十二の代数的ループは、理論物理学と生物学の歴史において前例のない構造を構成する。それらは、自然の基本定数——質量、結合定数、宇宙論的パラメータ、\(\pi\) の幾何、ゲノム構造、社会の動態、そして今や素数の分布——が、経験的に決定された数のランダムな集合ではなく、\textbf{単一の自己無撞着な代数的骨格の剛性的な射影} であることを示す。この骨格はたった三つの有理数によって定義される：\(S_{\mathrm{odd}} = 6/5 = 1.2\), \(S_{\mathrm{even}} = 4/5 = 0.8\)、およびそれらの比 \(3/2\)。
	
	この発見は、\textbf{ピタゴラス–プラトン的物理・生物現実観} への深い回帰を示すが、決定的な違いがある：それは神秘的な数秘術や哲学的思弁に基づくのではなく、厳密で反証可能な数学的枠組みに基づいている。宇宙は数の言語で書かれているという古い信念は、ここでその最初の具体的で科学的に検証可能な実現を、クォークから意識、安全な情報システム設計に至るまでの存在の全スペクトルにわたって見出す。
	
	\paragraph{避けられない科学的抵抗。}
	一握りの有理数が全ての観測可能な現実を支えているという主張は、科学コミュニティから深い懐疑をもって迎えられるべきであり、実際そうなるだろう。この懐疑は健全で必要である。証明の責任はこの枠組みに明確にある。しかし、本件のユニークさはその \textbf{数学的剛性と学際的範囲} にある。電子の質量を決定する同じ定数が、ヒトゲノムの構造も支配し、\(\pi\) を十万分の一まで近似し、二要素認証の効率を説明し、今や素数のゆらぎを予測する。このような多ループの一致が偶然に生じる確率は消えるほど小さい（\(p \ll 10^{-30}\)）。
	
	\section*{統一的反証基準：フレームワーク全体を危険にさらす}
	
	成熟した科学理論の定義的特徴は、反証されることをいとわないことである。YUCT 代数的枠組みは \textbf{剛性} である。調整可能な自由連続パラメータはない。二十二のループは以下の伝統的に独立した領域を接続する：
	\begin{itemize}
		\item \textbf{素粒子物理学：} フェルミオン質量と弱スケール（LHC で検証済み）。
		\item \textbf{重力物理学：} ブラックホールリンギングスケーリング（LIGO/Virgo/KAGRA）。
		\item \textbf{宇宙論：} 宇宙のバリオン質量（Planck, DESI）。
		\item \textbf{基礎数学：} \(\varphi\) と \(S_{\mathrm{odd}}\) からの \(\pi\) の創発；素数分布。
		\item \textbf{量子重力：} プランクスケールでの時間量子化。
		\item \textbf{分子生物学：} ゲノム構造、突然変異率、コドン使用（ゲノムデータベース）。
		\item \textbf{進化生物学：} 種分化率と双曲的成長（古生物学記録）。
		\item \textbf{社会学：} 臨界集団サイズと経済変動（歴史データ）。
		\item \textbf{情報理論とセキュリティ：} 2FA や量子もつれのような協調プロトコルの効率。
		\item \textbf{カオス理論：} ファイゲンバウム定数と秩序–カオス遷移。
	\end{itemize}
	いずれかの \textbf{ハードループ}（I–VI, XIV, XVI）における単一の明確で高信頼度の実験的失敗は、\textbf{代数フレームワーク全体を粉砕} するだろう。例えば：
	\begin{enumerate}
		\item 新しいフェルミオン質量が半整数ラダーから大きく逸脱するなら、ループ II を破壊する。
		\item リンギングスケーリング指数が明確に \(0.67\) でないなら、ループ VI を破壊する。
		\item 異なるゲノムにおける二塩基比が 1.5 から系統的に逸脱するなら、ループ XIV を破壊する。
		\item 双曲的成長が特異点時間窓を予測できないなら、ループ XII（ソフトだがその形式はハード）を破壊する。
	\end{enumerate}
	そのような失敗はいずれもパラメータ調整によって修復できない。なぜならパラメータがないからである。理論は単に間違っていることになる。これが代数的枠組みの究極の力である：それは基本法則の厳しく脆い美しさを提供し、科学コミュニティにその破壊を試みるよう招く。もしそれが生き残れば、現実の真の骨格としての地位を獲得するだろう。
	
	\section*{生物学的検証：トポロジカルシフト \(\sigma = 0.20\) からの細胞膜厚さ}
	\label{sec:membrane}
	
	YUCT 代数的枠組みの最も印象的な確認の一つは、基礎物理学からは程遠い領域、すなわち生きた細胞の構造からもたらされる。脂質二重層は、細胞の内部協調コア（DNA、タンパク質）を環境の熱雑音から分離する物理的境界である。その厚さは任意ではなく、ハドロン共鳴とフェルミオン質量ラダーを支配する同じトポロジカル不変量 \(\sigma = 0.20\) によって決定される。
	
	\subsection*{生細胞の協調要件}
	
	YUCT 生物学モジュール（付録~E）では、生細胞は情報漏洩を防ぐために最小効率 \(K_{\mathrm{eff}} > K_{\mathrm{eff,\,crit}} \approx 1.837\) を維持しなければならない孤立した協調核として扱われる。脂質膜は二つの情報相の間の界面として機能する。その幾何は同時に：
	\begin{itemize}
		\item 制御されないイオン拡散（エントロピー漏洩）をブロックするのに十分な厚さでなければならない。
		\item 迅速なシグナル伝達と触媒交換を可能にするのに十分な薄さでなければならない。
	\end{itemize}
	普遍的トポロジカルシフト \(\sigma\) は、両方の条件を保証する幾何学的スペーサーとして現れる。
	
	\subsection*{三項幾何からの膜厚さの導出（修正版）}
	
	三項空間 \(\mathbf{D} + \mathbf{I}\!\cdot\!\mathbf{R}\) において、偶数ループ定数 \(S_{\mathrm{even}} = 0.8\) は交替協調強度を定義する。シフト \(\sigma = 0.20\) は有効情報体積を圧縮する動径ギャップとして作用する。単一脂質リーフレットの疎水性コアについて、有効厚さ \(L_{\mathrm{mem}}\) は、基本スケーリング因子 \(3/2 = S_{\mathrm{odd}}/S_{\mathrm{even}}\) を単一アシル鎖長に適用し、二重層パッキングの立体障害を考慮するために \textbf{偶数セクター} のトポロジカルギャップを指数から差し引くことによって得られる：
	
	\[
	L_{\mathrm{mem}} = d_{\mathrm{lipid}} \cdot \left( \frac{3}{2} \right)^{S_{\mathrm{even}} - \sigma}
	= d_{\mathrm{lipid}} \cdot \left( \frac{3}{2} \right)^{0.8 - 0.20}
	= d_{\mathrm{lipid}} \cdot \left( \frac{3}{2} \right)^{0.6}.
	\]
	
	ここで \(d_{\mathrm{lipid}}\) は典型的な膜リン脂質（例えばパルミチン酸、\(d_{\mathrm{lipid}} \approx 2.5\)–\(3.0\)~nm）の完全に伸長した炭化水素尾部の長さである。指数 \(S_{\mathrm{even}} - \sigma = 0.6\) は、交替協調（偶数レベル）がパッキング幾何を支配することを反映し、シフト \(\sigma\) は実効指数を減少させ、奇数セクターのスケーリングを素朴に適用した場合に生じる非物理的伸長を防ぐ。
	
	二重層の総厚さはリーフレット厚さの二倍から、膜の流動性と熱的ゆらぎを考慮してゆらぎ振幅の二乗 \(\sigma^2\) を差し引いた小さな曲率補正を引いたものである：
	
	\[
	\text{膜厚さ} = 2 \, L_{\mathrm{mem}} \, (1 - \sigma^2)
	= 2 \cdot d_{\mathrm{lipid}} \cdot (3/2)^{0.6} \cdot (1 - 0.04).
	\]
	
	数値的には \((3/2)^{0.6} = e^{0.6 \ln 1.5} = e^{0.6 \times 0.405465} = e^{0.243279} \approx 1.2754\)。したがって：
	\[
	L_{\mathrm{mem}} \approx 1.2754 \; d_{\mathrm{lipid}}, \qquad
	\text{厚さ} = 2 \times 1.2754 \, d_{\mathrm{lipid}} \times 0.96 = 2.448 \, d_{\mathrm{lipid}}.
	\]
	
	\(d_{\mathrm{lipid}} = 3.0\)~nm（完全伸長パルミチン酸鎖長）を取ると：
	\[
	\text{厚さ} \approx 2.448 \times 3.0\ \text{nm} = 7.34\ \text{nm}.
	\]
	\(d_{\mathrm{lipid}} = 2.7\)~nm（混合脂質鎖の平均）を用いると、\(\approx 6.61\)~nm となり、実験範囲の下限にある。実際に測定されたヒト形質膜の厚さは \(7.5 \pm 0.5\)~nm であり、YUCT 予測を経験的窓の中心に置く。
	
	\subsection*{実験との比較と修正の必要性}
	
	ヒト細胞の形質膜の細胞学的測定は一貫して \textbf{7–8~nm} の厚さを与える（例えば、赤血球膜の電子顕微鏡を参照）。代数定数 \(S_{\mathrm{even}} = 0.8\), \(\sigma = 0.20\)、幾何因子 \(3/2\) から導出された修正 YUCT 予測は、正確にこの区間に収まる。以前の誤った公式（\(S_{\mathrm{odd}}-\sigma = 1.0\) を使用）は、厚さ \(3.12 \, d_{\mathrm{lipid}}\)（\(d_{\mathrm{lipid}}=3.0\)~nm なら約 9.4~nm）を与え、完全伸長鎖の最大可能二重層厚さを超え、立体障害制約に違反する。修正公式は物理的限界 \(L_{\mathrm{mem}} \le 2 d_{\mathrm{lipid}}\) を守り、三次元立体化学と両立する厚さを与える。
	
	\subsection*{代数的ループとの統合}
	
	この結果は \textbf{ソフトループ VIII}（代謝スケーリング）と \textbf{ループ XIV}（ゲノム構造）の直接の現れである。それは、生物学的アーキテクチャがランダム進化の産物だけではなく、粒子質量や宇宙論的パラメータを固定するのと同じ協調幾何によって制約されることを示す。細胞膜の厚さが定数 \(S_{\mathrm{even}} = 0.8\) と \(\sigma = 0.20\) から自由パラメータなしで計算できるという事実は、YUCT に強力な学際的テストを提供する。将来の測定が 7–8~nm の窓から有意に逸脱すれば、この特定の予測を反証し、トポロジカルシフトの普遍性に挑戦するだろう。
	
	\section*{データ利用可能性}
	全ての計算と補助導出は YPSDC リポジトリ \url{https://github.com/Alexey-Yakushev-YUCT/YPSDC} で入手可能である。
	
	\begin{thebibliography}{99}
		\bibitem{AppendixFerra} A. V. Yakushev, ``付録 Ferra1.2：秩序相と無秩序相の普遍協調定数,'' in \emph{YUCT}, Zenodo (2026).
		\bibitem{AppendixJ} A. V. Yakushev, ``付録 J：YUCT トポロジカルオフセットからの標準模型フレーバーパラメータの完全導出,'' in \emph{YUCT}, Zenodo (2026).
		\bibitem{AppendixLambda} A. V. Yakushev, ``付録 \(\Lambda\)：YUCT における階層問題と宇宙定数問題の解決,'' in \emph{YUCT}, Zenodo (2026).
		\bibitem{AppendixEhrenfest} A. V. Yakushev, ``付録 Ehrenfest：回転円盤パラドックスの協調的解釈と非ユークリッド幾何の創発,'' in \emph{YUCT}, Zenodo (2026).
		\bibitem{AppendixOmega} A. V. Yakushev, ``付録 \(\Omega\)：宇宙の全球回転と双極子回転半径からの新しい最小年齢,'' in \emph{YUCT}, Zenodo (2026).
		\bibitem{AppendixY} A. V. Yakushev, ``付録 Y：YUCT の数学的基礎——第一原理からの導出,'' in \emph{YUCT}, Zenodo (2026).
		\bibitem{AppendixV} A. V. Yakushev, ``付録 V：協調プロセスのエントロピーとしての時間,'' in \emph{YUCT}, Zenodo (2026).
		\bibitem{AppendixM} A. V. Yakushev, ``付録 M：YUCT 宇宙論——インフレーションを協調相転移として,'' in \emph{YUCT}, Zenodo (2026).
		\bibitem{AppendixE} A. V. Yakushev, ``付録 E：協調遺伝学——分子生物学における YUCT 原理,'' in \emph{YUCT}, Zenodo (2026).
		\bibitem{AppendixX} A. V. Yakushev, ``付録 X：YUCT とシャノン情報理論の一般化,'' in \emph{YUCT}, Zenodo (2026).
		\bibitem{AppendixPrimeN} A. V. Yakushev, ``付録 PrimeN：YUCT 素数協調ラダー,'' in \emph{YUCT}, Zenodo (2026).
		\bibitem{Planck2018} Planck 共同研究, \emph{Astron. Astrophys.} \textbf{641}, A6 (2020).
		\bibitem{UnityReality} A. V. Yakushev, ``YUCT 現実の統一：協調秩序 (\(\beta = 2/3\)) からファイゲンバウムカオス (\(\delta = 4.6692\)) へ,'' Zenodo (2026).
		\bibitem{Feigenbaum1978} M. J. Feigenbaum, ``Quantitative Universality for a Class of Nonlinear Transformations,'' J. Stat. Phys. \textbf{19}, 25–52 (1978).
	\end{thebibliography}
	
\end{document}